\begin{abstract}
Quantitative research requires choosing how an economic question should be
measured, estimated, implemented, and evaluated, not merely producing an
executable answer.  We study whether an autonomous quantitative researcher can
improve this capacity through its own research experience.  Our proposed
method, \method{}, asks whether economic structure and statistical reasoning
can guide informative observations about limitations in the system's working
methods and drive persistent harness revisions whose usefulness is tested
through subsequent research.  A \Worker{} executes tasks and an \Evolver{}
revises the complete persistent \Worker{} harness --- including prompts,
executable tools, skills, memory, middleware, validators, and control flow ---
while foundation-model configurations and the revision policy remain fixed
within each declared experiment.

We demonstrate that research-method experimentation can improve matched
initial-versus-frozen performance on backtesting tasks.  In a complete
controlled comparison on 20~BacktestBench instances, the frozen evolved
harness scores 16/20 against the initial harness's 14/20, with two task gains
and zero regressions.  The \Evolver{} diagnosed that the initial \Worker{}
evaluated round-trip performance using full-window daily mark-to-market P\&L,
conflating unrealized and realized gains across multi-day holds; the promoted
harness revision scopes evaluation to window-only indicators and realized
round-trip PNL.  Four fresh-\Worker{} consumer checks confirm consequential
policy uptake in the promoted candidate, including one zero-score negative
that bounds the claim.  The revised harness remains prompt-level; executable
tool and validator components are the target of ongoing development.

QuantCodeEval results are mixed across the same evolved harness: nine paired
tasks give binary 2/9 on the closed twenty-task checkpoint, and a separate
Pro \Evolver{} fork yields QCE binary 2/4 to~3/4 with one regression and QF
4/4 ties.  These results indicate that the learned backtesting-measurement
adaptation does not transfer uniformly across benchmark types at this stage,
which is consistent with the conditional nature of research-method
improvement: an applicable operation must match the current task's quantity,
information convention, and estimation context.

The completed BacktestBench comparison establishes one instance of
initial-to-frozen improvement through research-method experimentation.
Stable multi-task benefit and executable harness capability beyond prompt
revision remain open evaluation targets.
\end{abstract}
